% Appendices scaffold (Phase 2): five appendices A–E.
%
% Apps A–D follow genre_conventions_software.md §(h)–(k). App E is
% project-specific (symbolic-pipeline outputs for surveyed PGT theories,
% with the hx↔ax structural identity noted alongside the propagating-PGT
% equations).
%
% Restructure changes from the previous Phase 2 scaffold:
%   - Single-point trace-channel ghost diagnosis moved from §4.5 into
%     App D `ValidationPathological` (methodology validation, not a
%     standalone result).
%   - New `StabilityGuard` named paragraph in App D documenting the
%     stability-guard / IC plane-wave-snap engineering work that was
%     previously invisible in the report.
%   - App E equation-block subsections renamed to be theory-agnostic
%     (`SymbolicReference`, `SymbolicMechanism`, `SymbolicHeadline`)
%     so they can be filled at drafting time with whichever theories
%     prove most informative — the eventual content depends on
%     campaigns still in flight.
%   - All labels switched to bare CamelCase per STYLE_GUIDE.md §(c).
%
% Note: \appendix is invoked in main.tex BEFORE this file is input, so
% \section{} commands here automatically receive appendix numbering.

% ============================================================================
\section{TIDAL pipeline architecture}%
\label{Architecture}
% ============================================================================
% Target ~1800 w + diagram. Per genre_conventions_software.md §(h).

\paragraph*{Overview}
% TODO: 3–4 sentences précis of the whole pipeline — inputs, outputs,
% what is automated. Open the appendix; do not repeat the §3 summary.

% Figure A.1 — TIDAL architecture diagram.
\begin{figure}[t]
  \centering
  % TODO: \documentclass[border=10pt]{standalone}
\usepackage{tidal-tikz-styles}

\begin{document}
\begin{tikzpicture}[
    node distance=0.4cm
]

% ===================================================================
% ROW 1: Main pipeline stages (left to right)
% ===================================================================

% --- Input ---
\node[data] (toml) {\texttt{theory.toml}};

% --- Derive ---
\node[stage=Pink, right=0.6cm of toml] (derive)
    {Symbolic derivation\\[2pt]\scriptsize Wolfram/xAct};

% --- JSON (narrow waist) ---
\node[data, right=0.6cm of derive, minimum width=1.8cm,
      draw=BlueLine, fill=BlueFill, fill opacity=0.35,
      font=\footnotesize\bfseries] (json)
    {\texttt{spec.json}};

% --- Simulate ---
\node[stage=Purple, right=0.6cm of json] (sim)
    {Numerical integration};

% --- Snapshots ---
\node[data, right=0.6cm of sim] (snaps)
    {Field\\[-1pt]\scriptsize snapshots};

% --- Measure ---
\node[stage=Green, right=0.6cm of snaps] (meas)
    {Post-processing};

% --- Results ---
\node[data, right=0.6cm of meas] (results)
    {Measurement\\[-1pt]\scriptsize results};

% ===================================================================
% ARROWS (main flow)
% ===================================================================
\draw[arr] (toml) -- (derive);
\draw[arr] (derive) -- (json);
\draw[arr] (json) -- (sim);
\draw[arr] (sim) -- (snaps);
\draw[arr] (snaps) -- (meas);
\draw[arr] (meas) -- (results);


% ===================================================================
% BRACE: Symbolic vs Numerical separation at JSON
% ===================================================================
\draw[decorate, decoration={brace, amplitude=6pt, raise=2pt},
      PinkLine, thick]
    (toml.west |- derive.north) -- (derive.north east)
    node[midway, above=10pt, font=\scriptsize\bfseries, text=PinkLine]
    {Symbolic (CAS)};

\draw[decorate, decoration={brace, amplitude=6pt, raise=2pt},
      PurpleLine, thick]
    (sim.north west) -- (meas.north east)
    node[midway, above=10pt, font=\scriptsize\bfseries, text=PurpleLine]
    {Numerical (PDE)};

% JSON label
\node[above=0.25cm of json, font=\scriptsize\bfseries, text=BlueLine]
    {portable contract};

\end{tikzpicture}
\end{document}
 (TikZ)
  % TODO caption: "TIDAL pipeline. The symbolic stage (\xAct/\Mathematica)
  % ingests a TOML Lagrangian, derives the Euler–Lagrange equations,
  % decomposes by component, and exports a JSON specification. The
  % numerical stage (\Python) dispatches to one of five solver backends
  % based on equation properties, producing simulation data and
  % measurements."
  \caption{TODO: pipeline diagram (symbolic --> JSON --> numerical).}
  \label{ArchitectureDiagram}
\end{figure}

\paragraph*{Symbolic stage}
% TODO: 4–6 sentences. \xAct/\xPert, Euler–Lagrange derivation, component
% decomposition. User-perspective only — no internal routine descriptions.

\paragraph*{JSON handoff}
% TODO: 4–6 sentences. JSON schema fields present (equations, fields,
% canonical block, hamiltonian terms, etc.); reference json_schema.tex
% if appropriate.

\paragraph*{Solver dispatch}
% TODO: 4–6 sentences. Auto-selection logic — modal first if conditions
% are met, otherwise IDA/CVODE/leapfrog/scipy fallback.

\paragraph*{Output}
% TODO: 4–6 sentences. Simulation data files, measurement types
% (energy, conversion, mixing, spectrum, etc.).

\paragraph*{Dependencies}
% TODO: explicit citations to each package: \xAct, \xPert, xTras,
% SUNDIALS/scikit-sundae, \NumPy, \SciPy, \Mathematica. Cite via
% \cite{} for each — bib audit must verify these keys exist.

% ============================================================================
\section{Symbolic stage}%
\label{SymbolicStage}
% ============================================================================
% Target ~1600 w. Per genre_conventions_software.md §(i).

\paragraph*{xAct environment}
% TODO: which packages are loaded; why \xPert is used for metric
% perturbations (not xTens alone); why xTras is needed for the
% Einstein–Cartan coupling.

\paragraph*{Lagrangian input from TOML}
% TODO: how a TOML config becomes a Wolfram session; representative
% user-session listing showing how to invoke `tidal derive`.

\paragraph*{Euler--Lagrange derivation}
% TODO: E-L velocity form (state, v, A0); not Legendre transform.

\paragraph*{Component decomposition}
% TODO: DecomposeToComponents, cross-field handling via additionalFields.
% Brief note on why component-level E-L is required.

\paragraph*{JSON export}
% TODO: which fields are written; how the canonical block is constructed.

\paragraph*{Failure modes and diagnostics}
% TODO: derivation timeouts, index inconsistency, malformed JSON; how
% each is diagnosed.

% ============================================================================
\section{Numerical stage}%
\label{NumericalStage}
% ============================================================================
% Target ~2000 w. Per genre_conventions_software.md §(j).

\paragraph*{Solver selection hierarchy}
% TODO: opener — modal --> IDA --> CVODE --> leapfrog --> scipy;
% auto-selection logic: which equation properties trigger each backend.

\subsection{Fourier modal solver}%
\label{ModalSolver}
% Target ~800 w — primary backend, longest treatment.

\paragraph*{Eigendecomposition}
% TODO: constant-coefficient case — exact matrix exponential. State what
% "machine precision" means in this context.

\paragraph*{Position-dependent coefficients}
% TODO: Krylov \texttt{expm\_multiply}.

\paragraph*{Constraint elimination}
% TODO: Fourier Schur complement.

\paragraph*{Mathematical foundations}
% TODO: cite \cite{higham2008functions} and \cite{molervanloan2003nineteen}
% for matrix-exponential methods.

% Figure C.1 — Fourier modal solver convergence. Referenced from App D.
\begin{figure}[t]
  \centering
  % TODO: \includegraphics{figC1_solver_convergence}
  % TODO caption: "Fourier modal solver convergence: relative error vs grid
  % resolution for Einstein–Maxwell at $\Bzero = [\text{value}]$,
  % $D = [\text{value}]$. The solver attains machine precision
  % ($\sim 10^{-14}$) at modest resolution, demonstrating the
  % spectral-precision behaviour expected for a constant-coefficient
  % system. This figure also satisfies the App D convergence requirement."
  \caption{TODO: solver convergence vs resolution.}
  \label{SolverConvergence}
\end{figure}

\subsection{Analytical Jacobian}%
\label{AnalyticalJacobian}
% Target ~400 w.

\paragraph*{Three-tier structure}
% TODO: dense (N <= 2K), sparse (2K < N <= 200K), GMRES (N > 200K). One
% sentence on when each is active.

\paragraph*{Speedup factors}
% TODO: quote the measured speedups (5.3x dense; 2.5x IDA sparse;
% 1.2–1.4x CVODE sparse) from solver_optimizations.tex.

\subsection{Other backends}%
\label{OtherBackends}
% Target ~200 w total.

\paragraph*{IDA}
% TODO: 3–4 sentences on DAE / constraint handling.

\paragraph*{CVODE}
% TODO: 3–4 sentences on adaptive ODE for waves.

\paragraph*{Leapfrog}
% TODO: 3–4 sentences on symplectic integration.

\paragraph*{scipy}
% TODO: 3–4 sentences on general-purpose fallback.

\paragraph*{JSON-to-solver interface}
% TODO: closing paragraph with two distinct moves —
% (i) coefficient-matrix EXTRACTION from the JSON spec (which fields are
% read, how they are interpreted); (ii) solver DISPATCH (how
% auto-selection decides which backend to call given the extracted
% properties). One sentence each.

% ============================================================================
\section{Validation suite}%
\label{Validation}
% ============================================================================
% Target ~1300 w (was 1200 — +100 because the §4.5 ghost diagnosis is
% now hosted here as methodology validation, and the StabilityGuard
% paragraph is new). Per genre_conventions_software.md §(k) Barker
% validation hierarchy.

\paragraph*{Analytic limits}
% TODO: massless Klein–Gordon, Maxwell. Compare TIDAL output to textbook
% result. Optionally one figure per case, or a summary figure with
% subpanels.

\paragraph*{Reference-model reproduction}
% TODO: Boccaletti formula validation. Most important validation result.
% Cross-ref \cref{BoccalettiValidation} (main body) and report the
% quantitative agreement here.

% Pathological-cases paragraph — now hosts the trace-channel ghost
% diagnosis. The §4.5 result was moved here per the structural review
% (a single-point ghost diagnosis is methodology validation, not a
% standalone §4 result).
\paragraph*{Pathological cases}
% TODO: ~250 w. Demonstrates that TIDAL correctly diagnoses ghost
% instability at a known unstable parameter point. This validates the
% diagnosis methodology (TIDAL produces an exponentially-growing
% time-domain signal at parameter points where the no-ghost condition is
% violated). State the no-ghost condition as a closed inequality (eqn
% \cref{NoGhostCondition}) citing \cite{sezgin1980ghostfree} and
% \cite{lin2019pgt}; report whether the inequality is satisfied at the
% studied parameter point; cross-ref \cref{TraceKineticOutput} (App E)
% for the kinetic matrix that exhibits the wrong-sign coefficient.
% Frame as methodology validation, not as a sector-wide claim.
\begin{align}
  % TODO: no-ghost inequality
  \label{NoGhostCondition}
\end{align}

% Convergence paragraph — same figure (SolverConvergence) used by App C.
\paragraph*{Convergence}
% TODO: convergence test for the Fourier modal solver. Cross-ref
% \cref{SolverConvergence} (App C — same figure serves App D). State
% the convergence rate (spectral / machine-precision for the modal
% solver).

% NEW — StabilityGuard paragraph documents the engineering work that
% was previously invisible in the report (per Agent 3 review I7).
% Stability guard architecture, VHV crash fix, IC plane-wave-snap.
\paragraph*{Stability guard for plane-wave initial conditions}
% TODO: ~150 w. Document the stability guard that prevents spurious
% exponential growth from off-grid initial-condition spectral leakage.
% Off-grid `--ic-wavevector` values cause cos(k·x) to leak amplitude
% onto every discrete Fourier mode, including modes with tachyonic
% eigenvalues in PGT systems with cross-coupling — this triggers
% spurious SimulationDivergedError. TIDAL auto-snaps the IC wavevector
% to the nearest discrete mode (clamped below Nyquist) to eliminate
% the leakage. Reference docs/tex/plane_wave_ic.tex; cite GitHub issue
% #263 for the original observation. Frame as a methodology
% contribution that protects survey results from a specific
% well-defined failure mode.

% Figure D.1 — Stage A heatmap.
\begin{figure}[t]
  \centering
  % TODO: \includegraphics{figD1_stageA_heatmap}
  % TODO caption: "Stage A heatmap: $\Pconv^\text{max}$ in the
  % $(m_A, \varepsilon)$ parameter plane for the [sector] at
  % $\Bzero = [\text{value}]$, $t_\text{end} = [\text{value}]$. Bounds
  % the amplification across the swept region; supports the
  % corresponding row of \cref{SectorSurveyTable} (§4.2)."
  \caption{TODO: parameter-plane heatmap supporting a sector-survey row.}
  \label{StageAHeatmap}
\end{figure}

% Figure D.2 — Stage B sweep.
\begin{figure}[t]
  \centering
  % TODO: \includegraphics{figD2_stageB_sweep}
  % TODO caption: "Stage B sweep: scatter of $\Pconv^\text{max}$ vs the
  % analytic baseline $\sin^2(\kappa\Bzero t/2)$ across [N] runs in the
  % swept torsion-coupling region. The tight diagonal supports the
  % corresponding row of \cref{SectorSurveyTable}."
  \caption{TODO: sweep scatter of $\Pconv$ vs analytic baseline.}
  \label{StageBSweep}
\end{figure}

% Figure D.3 — Channel comparison.
\begin{figure}[t]
  \centering
  % TODO: \includegraphics{figD3_channel_comparison}
  % TODO caption: "Channel comparison: time evolution of conversion
  % probability in the $h_\times \leftrightarrow \mathrm{ax}$ channel
  % (solid) and the trace channel (dashed). The two channels behave
  % qualitatively differently, supporting the channel-decomposition
  % framing of \cref{ChannelStructure}."
  \caption{TODO: channel comparison; hx<->ax vs trace time series.}
  \label{ChannelComparison}
\end{figure}

% ============================================================================
\section{Symbolic-pipeline outputs for the surveyed theories}%
\label{SymbolicOutputs}
% ============================================================================
% Target ~700 w + long equations. Project-specific physics appendix.
%
% Subsection slots are intentionally theory-agnostic at scaffold time
% (`SymbolicReference`, `SymbolicMechanism`, `SymbolicHeadline`) so they
% can be filled with whichever theories prove most informative — the
% eventual choice depends on campaigns still in flight. The HxAxIdentity
% equation is hosted in `SymbolicMechanism` because the channel-structure
% identity is the most stable cross-reference target across the
% restructure.

\paragraph*{Purpose and reading guide}
% TODO: ~100 w paragraph. These are equations as produced by the TIDAL
% Wolfram pipeline for the PGT models surveyed in §4 that are too lengthy
% to display in the main body. Each block is the equation of motion or
% kinetic matrix in a chosen basis. Readers reproducing a §4 result can
% copy the equations into a Mathematica session and continue from there.
% Each block names the source JSON path so the equations can be
% regenerated from the pipeline.

\subsection{Reference theory: minimally-coupled Einstein--Cartan}%
\label{SymbolicReference}
% Target ~150 w + equation block. Source JSON: examples/<reference
% theory>/data/<file>.json. Serves as the simplest case in the survey
% — the theory whose kinetic matrix the reader should examine first to
% calibrate expectations for the other equation blocks.

\paragraph*{Kinetic matrix}
% TODO: introductory sentence + equation block. For long output, use
% \lstinputlisting{...} from an external .tex file
% (genre_conventions_software.md §(f.1)).
\begin{align}
  % TODO: paste from the relevant JSON. Source: examples/<reference>/data/...
  \label{ReferenceKineticOutput}
\end{align}

\paragraph*{Structural baseline}
% TODO: one sentence noting any structural feature that calibrates
% comparison with the other equation blocks (e.g. which entries are
% zero in the reference theory; which are non-zero only off-diagonal).

\subsection{Channel-structure mechanism: propagating PGT}%
\label{SymbolicMechanism}
% Target ~300 w + equation blocks. Source JSON: examples/<propagating
% PGT>/data/<file>.json. This subsection hosts the hx↔ax structural
% identity equation \cref{HxAxIdentity}, which is the cross-reference
% target for §4.3 and §2.3 (channel decomposition).

\paragraph*{Kinetic matrix and equations of motion}
% TODO: introductory sentence + long equation block.
\begin{align}
  % TODO: paste the kinetic matrix and EOMs from the propagating-PGT JSON.
  \label{PropagatingKineticOutput}
\end{align}

% Single named paragraph hosting the hx↔ax structural identity. This
% is inspection of two coefficients in the kinetic matrix — not a
% multi-step derivation — but the resulting identity is structural
% (visible in the symbolic output of every theory in the surveyed
% admissibility class), not a parameter-point coincidence.
\paragraph*{Structural identity: $h_\times$ and ax kinetic coefficients}
% TODO: ~80 w. State that the two relevant kinetic-matrix entries
% \cref{PropagatingKineticOutput} are identical at linear order under
% the admissibility constraint of gr-qc/0001010; display the
% equal-coefficient line as \cref{HxAxIdentity}; note that this
% structural identity is the kinematic explanation for the
% \cref{ChannelStructure} (§4.3) finding and the survey nulls in
% \cref{SectorSurveyTable} (§4.2).
\begin{align}
  % TODO: equal-coefficient line showing hx-coefficient = ax-coefficient
  \label{HxAxIdentity}
\end{align}

\subsection{Headline-theory equation block}%
\label{SymbolicHeadline}
% Target ~250 w + equation block. Source JSON: TBD — selected at
% drafting time to match the §4.4 headline finding. If the headline
% turns out to be a non-minimal-coupling amplification corner, this
% subsection hosts that theory's kinetic matrix (showing the new
% vertex that breaks the channel-structure identity). If the headline
% is structural rather than amplification, this subsection may host
% additional kinematic identities or wrong-sign coefficients.
%
% Also serves as the cross-reference target for the App D
% `Pathological cases` paragraph if that paragraph references a
% trace-channel kinetic-matrix excerpt.

\paragraph*{Kinetic matrix}
% TODO: introductory sentence + equation block at the headline
% parameter point or theory class.
\begin{align}
  % TODO: paste from the relevant JSON.
  \label{HeadlineKineticOutput}
\end{align}

% Reused label for the trace-channel kinetic-matrix excerpt (referenced
% from App D `Pathological cases`). Keeps a stable cross-reference
% target even though the surrounding subsection content is flexible.
\paragraph*{Trace-channel excerpt for the pathological-cases test}
% TODO: ~60 w. If the headline is *not* the trace-channel ghost (likely),
% this paragraph hosts only the trace-channel kinetic excerpt that App D
% Pathological-cases needs to cross-reference. If the headline IS the
% trace-channel ghost, fold this paragraph into the headline equation
% block above. The label \cref{TraceKineticOutput} is the stable
% cross-reference target for App D regardless.
\begin{align}
  % TODO: paste from the trace-channel JSON.
  \label{TraceKineticOutput}
\end{align}

\paragraph*{Structural feature of the headline regime}
% TODO: one paragraph noting what is visible in the equations that
% explains the §4.4 finding (e.g. a new non-zero off-diagonal entry,
% a wrong-sign kinetic coefficient, an algebraic relation that fails).
